\chapter{エントロピー}
    \section{連続型確率分布の場合}
        \subsection{一様分布がエントロピーを最大化する}
            \begin{shadebox}
                $a,b\in\realNumbers,\;a<b$とする。
                $p$を$[a,b]$上の確率密度関数とする。
                $p$のエントロピーが最大となるのは$p$が一様分布の時かつその時に限る。
            \end{shadebox}
            \begin{proof}
                \quad\par
                変分法を用いる。
                $[a,b]$上の確率分布であって密度関数が存在するもの全体の集合を$\mathcal{F}$とする。
                $P\in\mathcal{F}$には固定端条件$P(a)=0,P(b)=1$が課される。
                $\mathcal{F}$上の汎関数$I$を次式で定義する。
                \[ I: P\in\mathcal{F} \mapsto \integrate{a}{b}{f(P(x))}{}{x} \quad \text{where} \quad f(x) = -x\log x \]
                $f$は上に凸であることを容易に示せる。
                よって\ref{汎関数の大域的最小点の十分条件}より$p$がEuler方程式の解であることと、$p$がエントロピーを最大化することは同値である。
                Euler方程式を変形すると$\log P'(x) + 1 = 0$となり、$P'(x) = \mathrm{const.}$である。
                これと条件$P(a)=0, P(b)=1$より$p(x) = 1/(b-a)$となる。
            \end{proof}
